\documentclass[11pt]{article}
\usepackage{fontspec,amsmath,amssymb,amsthm,hyperref}
\usepackage[margin=1in]{geometry}
\setlength{\emergencystretch}{3em}
\title{Gaussian velocity moments and the stochastic local step}
\author{LangevinLean mathematical exposition}
\date{}
\begin{document}\maketitle
\noindent Source module: \texttt{KineticLangevin/LpOneStep.lean}. Compile this proof document with XeLaTeX or LuaLaTeX. The accompanying HTML pairs the mathematical argument with the exact Lean statements and records the scope of the formalization.


This module proves the exact \(\sqrt p\) reference velocity norm from a standard Gaussian marginal, then combines the genuine L\ensuremath{{}^{2}} nonlinear contraction with the actual exponential-integrator error. Its hypotheses display the Gaussian reference marginal and concrete exact integral paths; stationarity proving that marginal is a separate theorem.


\section{Obtain the exact velocity moment from its Gaussian law}


The standard Gaussian vector is the pushforward of independent scalar \(N(0,1)\) coordinates by \texttt{WithLp.toLp 2}. Compose that pushforward with coordinate evaluation; the finite-product marginal formula reduces its law to \(N(0,1)\), since every other coordinate measure has mass one. If an L\ensuremath{{}^{2}} velocity random variable has this standard Gaussian vector law, composition of laws gives the scalar coordinate laws. The project's proved Gaussian-coordinate moment identity then yields \(\|V\|_{L^2}=\sqrt{\#\iota}\). Thus the reference moment is derived from the explicit distribution rather than supplied as a numerical norm assumption.



\noindent Lean declarations: \texttt{standardGaussian\_coordinate\_law}, \texttt{gaussian\_velocity\_Lp\_norm}.

\section{Prove the concrete stochastic one-step inequality}


Let \(A\) be the auxiliary exact L\ensuremath{{}^{2}} path starting at the numerical phase point, and \(R\) a synchronously driven exact reference. Their transformed L\ensuremath{{}^{2}} distance is at most its initial value \(e\), because \(e^{-ms/\gamma}\le1\) for nonnegative time. Velocity projection and the exact Gaussian reference moment therefore give \(\|A_s^v\|_{L^2}\le\sqrt p+e\). Apply the actual mild-equation local-error theorem with this velocity bound: the numerical exponential update differs from \(A_h\) by at most \(3Mh^2(\sqrt p+e)/4\). Insert \(A_h\) with the transformed Hilbert triangle inequality and contract its distance to \(R_h\). The result is \(e_{\mathrm{next}}\le e^{-mh/\gamma}e+\frac34Mh^2(\sqrt p+e)\). Because the ambient space is L\ensuremath{{}^{2}}, the Bochner-integral and Hilbert triangle estimates are precisely the stochastic Minkowski estimates required by the paper.



The proof is factored through a shared theorem with an explicit reference velocity radius \(C\ge0\). A uniform bound \(\|V_t^{\rm ref}\|_{L^2}\le C\) supplies \(C+e_k\) in the local defect and the terminal bias \(Mh\sqrt2C/m\). The standard-Gaussian statement is a thin specialization with the radius \(\sqrt p\) proved from its marginal law. This same shared theorem applies to particle temperature \(1/N\), using the separately proved thermal radius \(\sqrt{(d_\theta+Nd_x)/N}\); there is no duplicate exponential-integrator argument.



\noindent Lean declarations: \texttt{lp\_exponentialStep\_coupling\_bound\_of\_velocity\_norm}, \texttt{lp\_exponentialStep\_coupling\_bound}.

\end{document}
